% =========================================================
% Slides Beamer Metropolis 16:9 -- Chapter 6 (ARCH and GARCH)
% Time Series Analysis -- Aymen Ben Brik (aymen.benbrik@esprit.tn)
%
% Extends _archives/source-slides/Non stationnarite.tex sections 4-5
% (volatility clustering, ARCH(q), GARCH(p,q), Engle LM test).
% =========================================================
\documentclass[aspectratio=169]{beamer}

% =========================================================
% Beamer preamble — Time Series Analysis module
% Theme : Metropolis, aspect ratio 16:9
% Author : Aymen Ben Brik (aymen.benbrik@esprit.tn)
% Esprit School of Business
%
% Include from each chapterX/slides.tex via :
%   % =========================================================
% Beamer preamble — Time Series Analysis module
% Theme : Metropolis, aspect ratio 16:9
% Author : Aymen Ben Brik (aymen.benbrik@esprit.tn)
% Esprit School of Business
%
% Include from each chapterX/slides.tex via :
%   % =========================================================
% Beamer preamble — Time Series Analysis module
% Theme : Metropolis, aspect ratio 16:9
% Author : Aymen Ben Brik (aymen.benbrik@esprit.tn)
% Esprit School of Business
%
% Include from each chapterX/slides.tex via :
%   \input{../preamble/slides-preamble.tex}
% AFTER the line \documentclass[aspectratio=169]{beamer}.
% =========================================================

\usepackage[utf8]{inputenc}
\usepackage[T1]{fontenc}
\usepackage[english]{babel}
\usepackage{lmodern}
\usepackage{amsmath, amssymb, mathtools, bm}
\usepackage{graphicx}
\usepackage{booktabs}
\usepackage{array}
\usepackage{multirow}
\usepackage{colortbl}
\usepackage{xcolor}
\usepackage{tikz}
\usetikzlibrary{positioning, arrows.meta, calc, shapes, fit, backgrounds, matrix, decorations.pathreplacing}
\usepackage{listings}
\usepackage[most]{tcolorbox}

% --- Theme ---
\usetheme{metropolis}
\setbeamertemplate{navigation symbols}{}
\metroset{block=fill}

% --- Esprit colour palette ---
\definecolor{espritBlue}{RGB}{30, 60, 110}
\definecolor{espritAccent}{RGB}{200, 80, 30}
\setbeamercolor{progress bar}{fg=espritAccent}
\setbeamercolor{title separator}{fg=espritAccent}
\setbeamercolor{frametitle}{bg=espritBlue}

% --- Code listings (Python + R, with UTF-8 literate) ---
\definecolor{codebg}{RGB}{248,248,248}
\definecolor{codeframe}{RGB}{220,220,220}
\definecolor{keyword}{RGB}{0,82,155}
\definecolor{comment}{RGB}{88,110,117}
\definecolor{string}{RGB}{133,153,0}

\lstdefinestyle{pythonstyle}{
  language=Python,
  basicstyle=\ttfamily\scriptsize,
  keywordstyle=\color{keyword}\bfseries,
  commentstyle=\color{comment}\itshape,
  stringstyle=\color{string},
  backgroundcolor=\color{codebg},
  frame=single,
  rulecolor=\color{codeframe},
  frameround=tttt,
  breaklines=true,
  showstringspaces=false,
  tabsize=2,
  literate=
    {é}{{\'e}}1 {è}{{\`e}}1 {ê}{{\^e}}1 {à}{{\`a}}1 {â}{{\^a}}1
    {ç}{{\c{c}}}1 {²}{{$^2$}}1 {°}{{${}^\circ$}}1
    {α}{{$\alpha$}}1 {β}{{$\beta$}}1 {γ}{{$\gamma$}}1
    {δ}{{$\delta$}}1 {θ}{{$\theta$}}1 {λ}{{$\lambda$}}1
    {μ}{{$\mu$}}1 {σ}{{$\sigma$}}1 {ε}{{$\varepsilon$}}1
    {ρ}{{$\rho$}}1 {τ}{{$\tau$}}1 {φ}{{$\phi$}}1
    {≤}{{$\leq$}}1 {≥}{{$\geq$}}1 {≈}{{$\approx$}}1
}
\lstdefinestyle{Rstyle}{
  language=R,
  basicstyle=\ttfamily\scriptsize,
  keywordstyle=\color{keyword}\bfseries,
  commentstyle=\color{comment}\itshape,
  stringstyle=\color{string},
  backgroundcolor=\color{codebg},
  frame=single,
  rulecolor=\color{codeframe},
  frameround=tttt,
  breaklines=true,
  showstringspaces=false,
  tabsize=2
}
\lstset{style=pythonstyle}

% --- Common macros ---
\input{../preamble/macros.tex}

% --- tcolorbox boxes for slides ---
\definecolor{boxTh}{RGB}{120, 30, 30}
\definecolor{boxProp}{RGB}{30, 100, 60}
\definecolor{boxDef}{RGB}{30, 60, 110}
\definecolor{boxEx}{RGB}{180, 120, 0}

\newtcolorbox{infobox}[1][]{enhanced, colback=espritBlue!5, colframe=espritBlue,
  fonttitle=\bfseries\small,
  title={\ifstrempty{#1}{}{#1}},
  boxrule=0.4pt, arc=2pt, left=4pt, right=4pt, top=2pt, bottom=2pt}
\newtcolorbox{warnbox}[1][]{enhanced, colback=espritAccent!5, colframe=espritAccent,
  fonttitle=\bfseries\small,
  title={\ifstrempty{#1}{}{#1}},
  boxrule=0.4pt, arc=2pt, left=4pt, right=4pt, top=2pt, bottom=2pt}
\newtcolorbox{thbox}[1][]{enhanced, colback=boxTh!5, colframe=boxTh,
  fonttitle=\bfseries\small,
  title={Theorem\ifstrempty{#1}{}{ -- #1}},
  boxrule=0.5pt, arc=2pt, left=4pt, right=4pt, top=2pt, bottom=2pt}
\newtcolorbox{propbox}[1][]{enhanced, colback=boxProp!5, colframe=boxProp,
  fonttitle=\bfseries\small,
  title={Proposition\ifstrempty{#1}{}{ -- #1}},
  boxrule=0.5pt, arc=2pt, left=4pt, right=4pt, top=2pt, bottom=2pt}
\newtcolorbox{defbox}[1][]{enhanced, colback=boxDef!5, colframe=boxDef,
  fonttitle=\bfseries\small,
  title={Definition\ifstrempty{#1}{}{ -- #1}},
  boxrule=0.5pt, arc=2pt, left=4pt, right=4pt, top=2pt, bottom=2pt}
\newtcolorbox{exbox}[1][]{enhanced, colback=boxEx!5, colframe=boxEx,
  fonttitle=\bfseries\small,
  title={Example\ifstrempty{#1}{}{ -- #1}},
  boxrule=0.5pt, arc=2pt, left=4pt, right=4pt, top=2pt, bottom=2pt}

% AFTER the line \documentclass[aspectratio=169]{beamer}.
% =========================================================

\usepackage[utf8]{inputenc}
\usepackage[T1]{fontenc}
\usepackage[english]{babel}
\usepackage{lmodern}
\usepackage{amsmath, amssymb, mathtools, bm}
\usepackage{graphicx}
\usepackage{booktabs}
\usepackage{array}
\usepackage{multirow}
\usepackage{colortbl}
\usepackage{xcolor}
\usepackage{tikz}
\usetikzlibrary{positioning, arrows.meta, calc, shapes, fit, backgrounds, matrix, decorations.pathreplacing}
\usepackage{listings}
\usepackage[most]{tcolorbox}

% --- Theme ---
\usetheme{metropolis}
\setbeamertemplate{navigation symbols}{}
\metroset{block=fill}

% --- Esprit colour palette ---
\definecolor{espritBlue}{RGB}{30, 60, 110}
\definecolor{espritAccent}{RGB}{200, 80, 30}
\setbeamercolor{progress bar}{fg=espritAccent}
\setbeamercolor{title separator}{fg=espritAccent}
\setbeamercolor{frametitle}{bg=espritBlue}

% --- Code listings (Python + R, with UTF-8 literate) ---
\definecolor{codebg}{RGB}{248,248,248}
\definecolor{codeframe}{RGB}{220,220,220}
\definecolor{keyword}{RGB}{0,82,155}
\definecolor{comment}{RGB}{88,110,117}
\definecolor{string}{RGB}{133,153,0}

\lstdefinestyle{pythonstyle}{
  language=Python,
  basicstyle=\ttfamily\scriptsize,
  keywordstyle=\color{keyword}\bfseries,
  commentstyle=\color{comment}\itshape,
  stringstyle=\color{string},
  backgroundcolor=\color{codebg},
  frame=single,
  rulecolor=\color{codeframe},
  frameround=tttt,
  breaklines=true,
  showstringspaces=false,
  tabsize=2,
  literate=
    {é}{{\'e}}1 {è}{{\`e}}1 {ê}{{\^e}}1 {à}{{\`a}}1 {â}{{\^a}}1
    {ç}{{\c{c}}}1 {²}{{$^2$}}1 {°}{{${}^\circ$}}1
    {α}{{$\alpha$}}1 {β}{{$\beta$}}1 {γ}{{$\gamma$}}1
    {δ}{{$\delta$}}1 {θ}{{$\theta$}}1 {λ}{{$\lambda$}}1
    {μ}{{$\mu$}}1 {σ}{{$\sigma$}}1 {ε}{{$\varepsilon$}}1
    {ρ}{{$\rho$}}1 {τ}{{$\tau$}}1 {φ}{{$\phi$}}1
    {≤}{{$\leq$}}1 {≥}{{$\geq$}}1 {≈}{{$\approx$}}1
}
\lstdefinestyle{Rstyle}{
  language=R,
  basicstyle=\ttfamily\scriptsize,
  keywordstyle=\color{keyword}\bfseries,
  commentstyle=\color{comment}\itshape,
  stringstyle=\color{string},
  backgroundcolor=\color{codebg},
  frame=single,
  rulecolor=\color{codeframe},
  frameround=tttt,
  breaklines=true,
  showstringspaces=false,
  tabsize=2
}
\lstset{style=pythonstyle}

% --- Common macros ---
% =========================================================
% Common math macros — Time Series Analysis module
% Author : Aymen Ben Brik (aymen.benbrik@esprit.tn)
% Esprit School of Business
% =========================================================

% --- Standard sets ---
\providecommand{\R}{\mathbb{R}}
\providecommand{\N}{\mathbb{N}}
\providecommand{\Z}{\mathbb{Z}}
\providecommand{\Q}{\mathbb{Q}}
\providecommand{\C}{\mathbb{C}}

% --- Probability / statistics ---
\providecommand{\E}{\mathbb{E}}
\providecommand{\Prob}{\mathbb{P}}
\providecommand{\Var}{\mathrm{Var}}
\providecommand{\Cov}{\mathrm{Cov}}
\providecommand{\Corr}{\mathrm{Corr}}
\providecommand{\indep}{\perp\!\!\!\perp}
\providecommand{\iid}{\overset{\text{i.i.d.}}{\sim}}
\providecommand{\given}{\,\big\vert\,}

% --- Estimators / hat ---
\providecommand{\hatm}{\widehat{m}}
\providecommand{\hatsig}{\widehat{\sigma}}
\providecommand{\hatrho}{\widehat{\rho}}
\providecommand{\hatphi}{\widehat{\phi}}
\providecommand{\hattheta}{\widehat{\theta}}

% --- Time series notation ---
\providecommand{\Lop}{L}                      % lag operator
\providecommand{\Diff}{\Delta}                % difference operator
\providecommand{\AR}[1]{\mathrm{AR}(#1)}
\providecommand{\MAm}[1]{\mathrm{MA}(#1)}
\providecommand{\ARMA}[2]{\mathrm{ARMA}(#1,#2)}
\providecommand{\ARIMA}[3]{\mathrm{ARIMA}(#1,#2,#3)}
\providecommand{\SARIMA}[6]{\mathrm{SARIMA}(#1,#2,#3)(#4,#5,#6)}
\providecommand{\ARCH}[1]{\mathrm{ARCH}(#1)}
\providecommand{\GARCH}[2]{\mathrm{GARCH}(#1,#2)}
\providecommand{\WN}{\mathrm{WN}}              % white noise
\providecommand{\MSE}{\mathrm{MSE}}
\providecommand{\MAE}{\mathrm{MAE}}
\providecommand{\RMSE}{\mathrm{RMSE}}
\providecommand{\AIC}{\mathrm{AIC}}
\providecommand{\BIC}{\mathrm{BIC}}
\providecommand{\ACF}{\mathrm{ACF}}
\providecommand{\PACF}{\mathrm{PACF}}

% --- Linear algebra ---
\providecommand{\transp}[1]{#1^{\!\top}}
\providecommand{\norm}[1]{\left\lVert #1 \right\rVert}
\providecommand{\abs}[1]{\left\lvert #1 \right\rvert}
\providecommand{\inner}[2]{\langle #1, #2 \rangle}
\providecommand{\diag}{\mathrm{diag}}
\providecommand{\tr}{\mathrm{tr}}

% --- Bold vectors / matrices ---
\providecommand{\vx}{\bm{x}}
\providecommand{\vy}{\bm{y}}
\providecommand{\vz}{\bm{z}}
\providecommand{\vh}{\bm{h}}
\providecommand{\vu}{\bm{u}}
\providecommand{\vv}{\bm{v}}
\providecommand{\vw}{\bm{w}}
\providecommand{\vb}{\bm{b}}
\providecommand{\vW}{\bm{W}}
\providecommand{\vSigma}{\bm{\Sigma}}
\providecommand{\veps}{\bm{\varepsilon}}

% --- Author identity (reused on titlepages) ---
\providecommand{\auteurNom}{Aymen Ben Brik}
\providecommand{\auteurEmail}{aymen.benbrik@esprit.tn}
\providecommand{\auteurInstitution}{Esprit School of Business}
\providecommand{\auteurRole}{Module head}
\providecommand{\moduleNom}{Time Series Analysis}
\providecommand{\anneeUniversitaire}{2025--2026}


% --- tcolorbox boxes for slides ---
\definecolor{boxTh}{RGB}{120, 30, 30}
\definecolor{boxProp}{RGB}{30, 100, 60}
\definecolor{boxDef}{RGB}{30, 60, 110}
\definecolor{boxEx}{RGB}{180, 120, 0}

\newtcolorbox{infobox}[1][]{enhanced, colback=espritBlue!5, colframe=espritBlue,
  fonttitle=\bfseries\small,
  title={\ifstrempty{#1}{}{#1}},
  boxrule=0.4pt, arc=2pt, left=4pt, right=4pt, top=2pt, bottom=2pt}
\newtcolorbox{warnbox}[1][]{enhanced, colback=espritAccent!5, colframe=espritAccent,
  fonttitle=\bfseries\small,
  title={\ifstrempty{#1}{}{#1}},
  boxrule=0.4pt, arc=2pt, left=4pt, right=4pt, top=2pt, bottom=2pt}
\newtcolorbox{thbox}[1][]{enhanced, colback=boxTh!5, colframe=boxTh,
  fonttitle=\bfseries\small,
  title={Theorem\ifstrempty{#1}{}{ -- #1}},
  boxrule=0.5pt, arc=2pt, left=4pt, right=4pt, top=2pt, bottom=2pt}
\newtcolorbox{propbox}[1][]{enhanced, colback=boxProp!5, colframe=boxProp,
  fonttitle=\bfseries\small,
  title={Proposition\ifstrempty{#1}{}{ -- #1}},
  boxrule=0.5pt, arc=2pt, left=4pt, right=4pt, top=2pt, bottom=2pt}
\newtcolorbox{defbox}[1][]{enhanced, colback=boxDef!5, colframe=boxDef,
  fonttitle=\bfseries\small,
  title={Definition\ifstrempty{#1}{}{ -- #1}},
  boxrule=0.5pt, arc=2pt, left=4pt, right=4pt, top=2pt, bottom=2pt}
\newtcolorbox{exbox}[1][]{enhanced, colback=boxEx!5, colframe=boxEx,
  fonttitle=\bfseries\small,
  title={Example\ifstrempty{#1}{}{ -- #1}},
  boxrule=0.5pt, arc=2pt, left=4pt, right=4pt, top=2pt, bottom=2pt}

% AFTER the line \documentclass[aspectratio=169]{beamer}.
% =========================================================

\usepackage[utf8]{inputenc}
\usepackage[T1]{fontenc}
\usepackage[english]{babel}
\usepackage{lmodern}
\usepackage{amsmath, amssymb, mathtools, bm}
\usepackage{graphicx}
\usepackage{booktabs}
\usepackage{array}
\usepackage{multirow}
\usepackage{colortbl}
\usepackage{xcolor}
\usepackage{tikz}
\usetikzlibrary{positioning, arrows.meta, calc, shapes, fit, backgrounds, matrix, decorations.pathreplacing}
\usepackage{listings}
\usepackage[most]{tcolorbox}

% --- Theme ---
\usetheme{metropolis}
\setbeamertemplate{navigation symbols}{}
\metroset{block=fill}

% --- Esprit colour palette ---
\definecolor{espritBlue}{RGB}{30, 60, 110}
\definecolor{espritAccent}{RGB}{200, 80, 30}
\setbeamercolor{progress bar}{fg=espritAccent}
\setbeamercolor{title separator}{fg=espritAccent}
\setbeamercolor{frametitle}{bg=espritBlue}

% --- Code listings (Python + R, with UTF-8 literate) ---
\definecolor{codebg}{RGB}{248,248,248}
\definecolor{codeframe}{RGB}{220,220,220}
\definecolor{keyword}{RGB}{0,82,155}
\definecolor{comment}{RGB}{88,110,117}
\definecolor{string}{RGB}{133,153,0}

\lstdefinestyle{pythonstyle}{
  language=Python,
  basicstyle=\ttfamily\scriptsize,
  keywordstyle=\color{keyword}\bfseries,
  commentstyle=\color{comment}\itshape,
  stringstyle=\color{string},
  backgroundcolor=\color{codebg},
  frame=single,
  rulecolor=\color{codeframe},
  frameround=tttt,
  breaklines=true,
  showstringspaces=false,
  tabsize=2,
  literate=
    {é}{{\'e}}1 {è}{{\`e}}1 {ê}{{\^e}}1 {à}{{\`a}}1 {â}{{\^a}}1
    {ç}{{\c{c}}}1 {²}{{$^2$}}1 {°}{{${}^\circ$}}1
    {α}{{$\alpha$}}1 {β}{{$\beta$}}1 {γ}{{$\gamma$}}1
    {δ}{{$\delta$}}1 {θ}{{$\theta$}}1 {λ}{{$\lambda$}}1
    {μ}{{$\mu$}}1 {σ}{{$\sigma$}}1 {ε}{{$\varepsilon$}}1
    {ρ}{{$\rho$}}1 {τ}{{$\tau$}}1 {φ}{{$\phi$}}1
    {≤}{{$\leq$}}1 {≥}{{$\geq$}}1 {≈}{{$\approx$}}1
}
\lstdefinestyle{Rstyle}{
  language=R,
  basicstyle=\ttfamily\scriptsize,
  keywordstyle=\color{keyword}\bfseries,
  commentstyle=\color{comment}\itshape,
  stringstyle=\color{string},
  backgroundcolor=\color{codebg},
  frame=single,
  rulecolor=\color{codeframe},
  frameround=tttt,
  breaklines=true,
  showstringspaces=false,
  tabsize=2
}
\lstset{style=pythonstyle}

% --- Common macros ---
% =========================================================
% Common math macros — Time Series Analysis module
% Author : Aymen Ben Brik (aymen.benbrik@esprit.tn)
% Esprit School of Business
% =========================================================

% --- Standard sets ---
\providecommand{\R}{\mathbb{R}}
\providecommand{\N}{\mathbb{N}}
\providecommand{\Z}{\mathbb{Z}}
\providecommand{\Q}{\mathbb{Q}}
\providecommand{\C}{\mathbb{C}}

% --- Probability / statistics ---
\providecommand{\E}{\mathbb{E}}
\providecommand{\Prob}{\mathbb{P}}
\providecommand{\Var}{\mathrm{Var}}
\providecommand{\Cov}{\mathrm{Cov}}
\providecommand{\Corr}{\mathrm{Corr}}
\providecommand{\indep}{\perp\!\!\!\perp}
\providecommand{\iid}{\overset{\text{i.i.d.}}{\sim}}
\providecommand{\given}{\,\big\vert\,}

% --- Estimators / hat ---
\providecommand{\hatm}{\widehat{m}}
\providecommand{\hatsig}{\widehat{\sigma}}
\providecommand{\hatrho}{\widehat{\rho}}
\providecommand{\hatphi}{\widehat{\phi}}
\providecommand{\hattheta}{\widehat{\theta}}

% --- Time series notation ---
\providecommand{\Lop}{L}                      % lag operator
\providecommand{\Diff}{\Delta}                % difference operator
\providecommand{\AR}[1]{\mathrm{AR}(#1)}
\providecommand{\MAm}[1]{\mathrm{MA}(#1)}
\providecommand{\ARMA}[2]{\mathrm{ARMA}(#1,#2)}
\providecommand{\ARIMA}[3]{\mathrm{ARIMA}(#1,#2,#3)}
\providecommand{\SARIMA}[6]{\mathrm{SARIMA}(#1,#2,#3)(#4,#5,#6)}
\providecommand{\ARCH}[1]{\mathrm{ARCH}(#1)}
\providecommand{\GARCH}[2]{\mathrm{GARCH}(#1,#2)}
\providecommand{\WN}{\mathrm{WN}}              % white noise
\providecommand{\MSE}{\mathrm{MSE}}
\providecommand{\MAE}{\mathrm{MAE}}
\providecommand{\RMSE}{\mathrm{RMSE}}
\providecommand{\AIC}{\mathrm{AIC}}
\providecommand{\BIC}{\mathrm{BIC}}
\providecommand{\ACF}{\mathrm{ACF}}
\providecommand{\PACF}{\mathrm{PACF}}

% --- Linear algebra ---
\providecommand{\transp}[1]{#1^{\!\top}}
\providecommand{\norm}[1]{\left\lVert #1 \right\rVert}
\providecommand{\abs}[1]{\left\lvert #1 \right\rvert}
\providecommand{\inner}[2]{\langle #1, #2 \rangle}
\providecommand{\diag}{\mathrm{diag}}
\providecommand{\tr}{\mathrm{tr}}

% --- Bold vectors / matrices ---
\providecommand{\vx}{\bm{x}}
\providecommand{\vy}{\bm{y}}
\providecommand{\vz}{\bm{z}}
\providecommand{\vh}{\bm{h}}
\providecommand{\vu}{\bm{u}}
\providecommand{\vv}{\bm{v}}
\providecommand{\vw}{\bm{w}}
\providecommand{\vb}{\bm{b}}
\providecommand{\vW}{\bm{W}}
\providecommand{\vSigma}{\bm{\Sigma}}
\providecommand{\veps}{\bm{\varepsilon}}

% --- Author identity (reused on titlepages) ---
\providecommand{\auteurNom}{Aymen Ben Brik}
\providecommand{\auteurEmail}{aymen.benbrik@esprit.tn}
\providecommand{\auteurInstitution}{Esprit School of Business}
\providecommand{\auteurRole}{Module head}
\providecommand{\moduleNom}{Time Series Analysis}
\providecommand{\anneeUniversitaire}{2025--2026}


% --- tcolorbox boxes for slides ---
\definecolor{boxTh}{RGB}{120, 30, 30}
\definecolor{boxProp}{RGB}{30, 100, 60}
\definecolor{boxDef}{RGB}{30, 60, 110}
\definecolor{boxEx}{RGB}{180, 120, 0}

\newtcolorbox{infobox}[1][]{enhanced, colback=espritBlue!5, colframe=espritBlue,
  fonttitle=\bfseries\small,
  title={\ifstrempty{#1}{}{#1}},
  boxrule=0.4pt, arc=2pt, left=4pt, right=4pt, top=2pt, bottom=2pt}
\newtcolorbox{warnbox}[1][]{enhanced, colback=espritAccent!5, colframe=espritAccent,
  fonttitle=\bfseries\small,
  title={\ifstrempty{#1}{}{#1}},
  boxrule=0.4pt, arc=2pt, left=4pt, right=4pt, top=2pt, bottom=2pt}
\newtcolorbox{thbox}[1][]{enhanced, colback=boxTh!5, colframe=boxTh,
  fonttitle=\bfseries\small,
  title={Theorem\ifstrempty{#1}{}{ -- #1}},
  boxrule=0.5pt, arc=2pt, left=4pt, right=4pt, top=2pt, bottom=2pt}
\newtcolorbox{propbox}[1][]{enhanced, colback=boxProp!5, colframe=boxProp,
  fonttitle=\bfseries\small,
  title={Proposition\ifstrempty{#1}{}{ -- #1}},
  boxrule=0.5pt, arc=2pt, left=4pt, right=4pt, top=2pt, bottom=2pt}
\newtcolorbox{defbox}[1][]{enhanced, colback=boxDef!5, colframe=boxDef,
  fonttitle=\bfseries\small,
  title={Definition\ifstrempty{#1}{}{ -- #1}},
  boxrule=0.5pt, arc=2pt, left=4pt, right=4pt, top=2pt, bottom=2pt}
\newtcolorbox{exbox}[1][]{enhanced, colback=boxEx!5, colframe=boxEx,
  fonttitle=\bfseries\small,
  title={Example\ifstrempty{#1}{}{ -- #1}},
  boxrule=0.5pt, arc=2pt, left=4pt, right=4pt, top=2pt, bottom=2pt}


\title[\moduleNom\ -- Ch.6]{ARCH and GARCH Models}
\subtitle{Volatility clustering $\cdot$ Conditional heteroskedasticity $\cdot$ MLE $\cdot$ VaR}
\author{\auteurNom \\ \footnotesize\auteurEmail}
\institute{\auteurInstitution}
\date{\anneeUniversitaire}

\begin{document}

\begin{frame}\titlepage\end{frame}

\begin{frame}{Learning objectives}
  By the end of this chapter, you will be able to:
  \begin{itemize}
    \item recognise the limits of homoskedastic ARMA on financial
          returns;
    \item state the ARCH$(q)$ and GARCH$(p,q)$ models, their
          stationarity conditions and unconditional moments;
    \item run Engle's LM test for ARCH effects on the residuals of a
          mean model;
    \item estimate a GARCH model by maximum likelihood and validate
          the standardised residuals;
    \item produce a one-step-ahead conditional volatility forecast
          and a Value-at-Risk estimate.
  \end{itemize}
\end{frame}

\begin{frame}{Chapter outline}
  \begin{enumerate}
    \item Stylised facts of financial returns
    \item Conditional vs unconditional variance
    \item ARCH$(q)$ model
    \item GARCH$(p,q)$ model
    \item Estimation: maximum likelihood
    \item Diagnostics: Engle LM, standardised residuals
    \item Volatility forecasting and Value-at-Risk
    \item Extensions: EGARCH, GJR-GARCH (overview)
  \end{enumerate}
\end{frame}

% =====================================================================
\section{Stylised facts of financial returns}

\begin{frame}{1.~Why ARMA is not enough}
  Standard ARMA models assume a constant variance
  $\Var(\varepsilon_t) = \sigma^2$. On financial returns
  $r_t = \log(P_t / P_{t-1})$, this assumption fails:
  \begin{itemize}
    \item the level looks (almost) stationary, but
    \item large (positive or negative) returns cluster in time;
    \item small returns also cluster in time;
    \item the unconditional distribution has heavier tails than
          $\mathcal N(0, \sigma^2)$.
  \end{itemize}
  \pause
  \alert{Mandelbrot (1963):} \emph{``Large changes tend to be
  followed by large changes, of either sign, and small changes by
  small changes.''}
\end{frame}

\begin{frame}{Stylised facts (financial returns)}
  \begin{itemize}
    \item $\E[r_t] \approx 0$, often a tiny drift.
    \item Sample ACF of $r_t$: nearly flat (returns are weakly
          serially uncorrelated).
    \item Sample ACF of $r_t^2$ (or $|r_t|$):
          \emph{strong, slowly decaying} autocorrelation
          $\Rightarrow$ \alert{volatility clustering}.
    \item Heavy-tailed marginal distribution: kurtosis
          $\widehat\kappa > 3$ (excess kurtosis).
    \item Leverage effect: negative shocks raise volatility more
          than positive shocks of the same magnitude.
  \end{itemize}
\end{frame}

% =====================================================================
\section{Conditional vs unconditional variance}

\begin{frame}{2.~Conditional vs unconditional variance}
  Information set: $\mathcal F_{t-1} = \sigma(r_{t-1}, r_{t-2}, \dots)$.
  \begin{itemize}
    \item \emph{Unconditional} variance:
          $\sigma^2 = \Var(r_t)$ -- the long-run, average risk.
    \item \emph{Conditional} variance:
          $\sigma_t^2 = \Var(r_t \mid \mathcal F_{t-1})$ -- the
          short-run, predictable risk.
  \end{itemize}
  \pause
  ARCH/GARCH models specify \emph{$\sigma_t^2$} as a function of
  past information, while keeping the unconditional variance
  $\sigma^2$ finite. The series is then
  \emph{unconditionally homoskedastic} but
  \emph{conditionally heteroskedastic}.
\end{frame}

\begin{frame}{The two-equation structure}
  Most volatility models split the dynamics in two:
  \begin{block}{Mean equation}
    \[
      r_t = \mu_t + \varepsilon_t,
      \qquad
      \varepsilon_t = \sigma_t\, z_t,
      \qquad
      z_t \iid \mathcal D(0, 1).
    \]
    $\mu_t$ is, e.g., a constant or an ARMA model.
  \end{block}
  \begin{block}{Variance equation}
    A specification of $\sigma_t^2$ in terms of past
    $\varepsilon_{t-i}^2$ and past $\sigma_{t-j}^2$.
  \end{block}
  \pause
  Common innovation distributions $\mathcal D$: standard normal
  (Gaussian QMLE), Student-$t$ for fat tails, GED.
\end{frame}

% =====================================================================
\section{ARCH$(q)$ model}

\begin{frame}{3.~ARCH$(q)$, Engle (1982)}
  \begin{block}{Definition}
    \[
      \varepsilon_t = \sigma_t\, z_t,
      \qquad
      \sigma_t^2 = \alpha_0
                  + \sum_{i = 1}^{q} \alpha_i\, \varepsilon_{t-i}^2,
    \]
    with $z_t \iid \mathcal N(0, 1)$, $\alpha_0 > 0$,
    $\alpha_i \geq 0$ for $i = 1, \dots, q$.
  \end{block}
  \pause
  \begin{itemize}
    \item Conditional variance reacts to the \emph{magnitude} of
          past shocks (squares, sign-blind).
    \item Robert Engle received the 2003 Nobel Prize in Economics
          for this model.
  \end{itemize}
\end{frame}

\begin{frame}{Properties of ARCH$(q)$}
  Assume $\sum_{i=1}^{q} \alpha_i < 1$ (covariance-stationarity).
  Then:
  \[
    \E[\varepsilon_t] = 0,
    \qquad
    \E[\varepsilon_t^2] = \frac{\alpha_0}{1 - \sum_{i = 1}^{q} \alpha_i}
        \;=\; \sigma^2 \quad (\text{unconditional variance}).
  \]
  \pause
  \begin{itemize}
    \item $(\varepsilon_t)$ is white noise (serially uncorrelated)
          but \emph{not} independent.
    \item $(\varepsilon_t^2)$ follows an AR$(q)$ in disguise:
          $\varepsilon_t^2 = \alpha_0 + \sum \alpha_i
          \varepsilon_{t-i}^2 + \nu_t$ with $\nu_t = \varepsilon_t^2
          - \sigma_t^2$.
    \item Excess kurtosis $\widehat\kappa > 3$: ARCH naturally
          generates fat tails.
  \end{itemize}
\end{frame}

% =====================================================================
\section{GARCH$(p, q)$ model}

\begin{frame}{4.~GARCH$(p, q)$, Bollerslev (1986)}
  ARCH$(q)$ requires many parameters to capture long memory in
  volatility. The \emph{generalised ARCH} model adds lagged
  conditional variances:
  \begin{block}{Definition}
    \[
      \sigma_t^2 = \alpha_0
                  + \sum_{i = 1}^{q} \alpha_i\, \varepsilon_{t-i}^2
                  + \sum_{j = 1}^{p} \beta_j\, \sigma_{t-j}^2,
    \]
    $\alpha_0 > 0$, $\alpha_i, \beta_j \geq 0$.
  \end{block}
  \pause
  \alert{Workhorse:} GARCH$(1, 1)$ -- two parameters
  $(\alpha_1, \beta_1)$ on top of $\alpha_0$, captures most of the
  volatility persistence in financial data.
\end{frame}

\begin{frame}{GARCH$(1, 1)$: anatomy}
  \[
    \sigma_t^2
    = \underbrace{\alpha_0}_{\text{baseline}}
    + \underbrace{\alpha_1\, \varepsilon_{t-1}^2}_{\text{news / shock}}
    + \underbrace{\beta_1\, \sigma_{t-1}^2}_{\text{persistence / memory}}.
  \]
  \pause
  \begin{itemize}
    \item $\alpha_1$: how strongly today's volatility reacts to
          yesterday's surprise.
    \item $\beta_1$: how strongly today's volatility remembers
          yesterday's volatility.
    \item $\alpha_1 + \beta_1$ close to $1$: very persistent
          volatility (typical for daily equity returns).
  \end{itemize}
\end{frame}

\begin{frame}{Stationarity and unconditional variance}
  \begin{block}{Covariance-stationarity condition}
    GARCH$(1,1)$ has a finite unconditional variance if and only if
    \[
      \alpha_1 + \beta_1 < 1.
    \]
    The unconditional variance is then
    \[
      \sigma^2 = \E[\varepsilon_t^2]
              = \frac{\alpha_0}{1 - \alpha_1 - \beta_1}.
    \]
  \end{block}
  \pause
  General GARCH$(p, q)$: replace $\alpha_1 + \beta_1$ by
  $\sum \alpha_i + \sum \beta_j$.
  \alert{Warning.} If $\alpha_1 + \beta_1 = 1$ we have an
  \emph{IGARCH} -- shocks have permanent effects on volatility.
\end{frame}

% =====================================================================
\section{Estimation}

\begin{frame}{5.~MLE for GARCH}
  Conditional Gaussian density
  $r_t \mid \mathcal F_{t-1} \sim \mathcal N(\mu, \sigma_t^2)$:
  \[
    \log \mathcal L(\bm\theta)
    = -\tfrac{T}{2} \log(2 \pi)
      -\tfrac{1}{2} \sum_{t = 1}^{T}
          \Bigl(
              \log \sigma_t^2(\bm\theta)
              + \tfrac{(r_t - \mu)^2}{\sigma_t^2(\bm\theta)}
          \Bigr).
  \]
  \begin{itemize}
    \item $\sigma_t^2$ is computed recursively from the model
          equation.
    \item Initialised with $\sigma_1^2 = \widehat\sigma^2$ (sample
          variance) by convention.
    \item Maximisation by quasi-Newton; the Gaussian QMLE is
          consistent and asymptotically normal even if the true
          $z_t$ is non-Gaussian (heavy-tailed).
  \end{itemize}
\end{frame}

\begin{frame}[fragile]{Estimation -- Python (\texttt{arch} package)}
\begin{lstlisting}[language=Python]
from arch import arch_model

am  = arch_model(returns * 100,        # rescale for numerics
                 mean="Constant",
                 vol="GARCH",
                 p=1, q=1, dist="normal")
res = am.fit(disp="off")
print(res.summary())
print("alpha + beta =",
      res.params["alpha[1]"] + res.params["beta[1]"])
\end{lstlisting}
  \begin{itemize}
    \item \texttt{vol="ARCH"} for pure ARCH$(q)$.
    \item \texttt{dist="t"} for Student-$t$ innovations.
    \item \texttt{res.conditional\_volatility} stores
          $\widehat\sigma_t$.
  \end{itemize}
\end{frame}

% =====================================================================
\section{Diagnostics}

\begin{frame}{6.~Engle's LM test for ARCH effects}
  Run the auxiliary regression
  \[
    \widehat\varepsilon_t^2
    = \gamma_0 + \gamma_1 \widehat\varepsilon_{t-1}^2
              + \dots + \gamma_q \widehat\varepsilon_{t-q}^2 + u_t.
  \]
  \begin{itemize}
    \item $H_0$: $\gamma_1 = \dots = \gamma_q = 0$ (no ARCH effect).
    \item Test statistic: $LM = T \cdot R^2$, with $T$ the sample
          size and $R^2$ the regression's coefficient of
          determination.
    \item Under $H_0$, $LM \xrightarrow{d} \chi^2_q$.
    \item Reject $H_0$ if $LM > \chi^2_{q, 1 - \alpha}$ -- evidence
          for ARCH effects, fit a GARCH model.
  \end{itemize}
\end{frame}

\begin{frame}{Standardised residual diagnostics}
  After fitting GARCH, check
  \[
    \widehat z_t = \frac{\widehat\varepsilon_t}{\widehat\sigma_t}.
  \]
  A well-fitted model satisfies:
  \begin{itemize}
    \item Sample ACF of $\widehat z_t$ inside Bartlett bands.
    \item Sample ACF of $\widehat z_t^2$ also inside the bands
          (volatility clustering removed).
    \item Engle LM on $\widehat z_t^2$ does \emph{not} reject.
    \item Marginal distribution closer to the assumed
          $\mathcal D(0, 1)$ (QQ-plot).
  \end{itemize}
\end{frame}

% =====================================================================
\section{Forecasting and Value-at-Risk}

\begin{frame}{7.~One-step-ahead variance forecast}
  For GARCH$(1, 1)$,
  \[
    \widehat\sigma_{T+1}^2
    = \alpha_0 + \alpha_1 \widehat\varepsilon_T^2
               + \beta_1 \widehat\sigma_T^2.
  \]
  Multi-step iterated forecast (mean-reverting to the
  unconditional variance):
  \[
    \widehat\sigma_{T+h}^2
    = \sigma^2
    + (\alpha_1 + \beta_1)^{h - 1}
        \bigl(\widehat\sigma_{T+1}^2 - \sigma^2\bigr).
  \]
  \pause
  $\widehat\sigma_{T+h}^2 \to \sigma^2$ as $h \to \infty$, at a rate
  controlled by the persistence $\alpha_1 + \beta_1$.
\end{frame}

\begin{frame}{Value-at-Risk}
  At horizon $1$, level $\alpha$ (e.g.\ $1\%$), assuming Gaussian
  innovations:
  \[
    \mathrm{VaR}_{T+1}(\alpha)
    = -\bigl(\widehat\mu + z_\alpha\, \widehat\sigma_{T+1}\bigr),
  \]
  with $z_\alpha = \Phi^{-1}(\alpha)$ -- e.g.\
  $z_{0.01} = -2.326$, $z_{0.05} = -1.645$.
  \pause
  \begin{itemize}
    \item Static (unconditional) VaR uses sample $\sigma$ instead of
          $\widehat\sigma_{T+1}$.
    \item GARCH-VaR is responsive to the current volatility regime
          and produces tighter intervals in calm periods, wider
          ones in turbulent periods.
  \end{itemize}
\end{frame}

% =====================================================================
\section{Extensions (overview)}

\begin{frame}{8.~Extensions}
  \begin{block}{EGARCH (Nelson, 1991)}
    Models $\log \sigma_t^2$ to relax non-negativity constraints
    and capture asymmetric responses (leverage effect):
    \[
      \log \sigma_t^2 = \omega + \alpha\, g(z_{t-1}) + \beta \log
      \sigma_{t-1}^2.
    \]
  \end{block}
  \begin{block}{GJR-GARCH (Glosten--Jagannathan--Runkle, 1993)}
    Adds a threshold term that activates only on negative shocks:
    \[
      \sigma_t^2 = \omega + (\alpha + \gamma\, \mathbb 1_{\varepsilon_{t-1}<0})
                            \varepsilon_{t-1}^2
                          + \beta\, \sigma_{t-1}^2.
    \]
    $\gamma > 0$: bad news raises volatility more than good news.
  \end{block}
\end{frame}

\begin{frame}{Take-home messages}
  \begin{itemize}
    \item Returns look uncorrelated, but their squares are
          \emph{strongly} autocorrelated -- volatility clustering.
    \item ARCH$(q)$ models conditional variance as an AR on past
          squared shocks; GARCH$(p, q)$ adds memory via past
          variances.
    \item GARCH$(1, 1)$ is the standard workhorse:
          $\alpha_1 + \beta_1 < 1$ for stationarity, with
          unconditional variance $\alpha_0 / (1 - \alpha_1 - \beta_1)$.
    \item Test for ARCH effects with Engle's LM; estimate by MLE
          (\texttt{arch} package); validate with standardised
          residuals.
    \item Conditional VaR uses $\widehat\sigma_{T+1}$ instead of
          the sample $\widehat\sigma$ -- it adapts to the current
          regime.
  \end{itemize}
\end{frame}

\end{document}
